% ==================== قسمت (الف) ====================
\subsectionaddtolist{الف}

پالس مربعی متناوب $p(t)$ با دوره‌ی $T$، دامنه‌ی $1$ و عرض $d$ را در نظر می‌گیریم:
$$p(t) = \begin{cases} 1 & 0 \leq t < d \\ 0 & d \leq t < T \end{cases}, \qquad p(t+T)=p(t)$$

ضرایب سری فوریه‌ی $p(t)$:
$$p_n = \frac{1}{T}\int_0^d e^{-j\frac{2\pi n}{T}t}\,dt$$

 برای $n=0$:
  \[
  p_0 = \dfrac{d}{T}
  \]

 برای $n\neq 0$: 
    
 \[
 \displaystyle p_n = \frac{1-e^{-j\frac{2\pi nd}{T}}}{j2\pi n}
 \]

مشتق توزیعی $p(t)$ یک قطار ضربه‌ی ترکیبی است:
$$p'(t) = \sum_{k=-\infty}^{\infty}\bigl[\delta(t-kT) - \delta(t-d-kT)\bigr]$$

ضرایب سری فوریه‌ی $p'(t)$ با انتگرال مستقیم:
$$[p']_n = \frac{1}{T}\int_0^T p'(t)\,e^{-j\frac{2\pi n}{T}t}\,dt
= \frac{1}{T}\bigl(e^{0} - e^{-j\frac{2\pi nd}{T}}\bigr)
= \frac{1-e^{-j\frac{2\pi nd}{T}}}{T}$$

حال رابطه‌ی صورت سوال را بررسی می‌کنیم:

$$j\frac{2\pi n}{T}\,p_n = j\frac{2\pi n}{T} \cdot \frac{1-e^{-j\frac{2\pi nd}{T}}}{j2\pi n}
= \frac{1-e^{-j\frac{2\pi nd}{T}}}{T} = [p']_n $$

برای $n=0$: هر دو طرف برابر صفرند چون میانگین مشتق یک سیگنال متناوب صفر است.

پس رابطه‌ی $[p']_n = j\dfrac{2\pi n}{T}p_n$ برای پالس مربعی و قطار ضربه‌ی مشتق آن برقرار است.

% ==================== قسمت (ب) ====================
\subsectionaddtolist{ب}

مشتق توزیعی $f(t)$ برابر است با:
$$f'(t) = f'_{\text{smooth}}(t) + \sum_i \Delta_i\,\delta(t-t_i)$$

ضرایب فوریه را با انتگرال‌گیری روی قطعه‌های پیوسته و اعمال انتگرال‌گیری جزء‌به‌جزء محاسبه می‌کنیم. با تقسیم دوره‌ی $[0,T)$ به قطعه‌های پیوسته و اعمال انتگرال‌گیری جزء‌به‌جزء روی هر قطعه:

\begin{align*}
[f']_n &= \frac{1}{T}\int_0^T f'(t)\,e^{-j\frac{2\pi n}{T}t}\,dt \\
&= \frac{1}{T}\left[\,\Bigl[f(t)e^{-j\frac{2\pi n}{T}t}\Bigr]_0^T
   - \sum_i \Delta_i\,e^{-j\frac{2\pi n t_i}{T}}
   + j\frac{2\pi n}{T}\int_0^T f(t)\,e^{-j\frac{2\pi n}{T}t}\,dt\right]
   + \frac{1}{T}\sum_i \Delta_i\,e^{-j\frac{2\pi n t_i}{T}}
\end{align*}

جمله‌ی مرزی به دلیل تناوب از بین می‌رود ($f(T)e^{-j2\pi n}=f(0)$)، و جملات $\Delta_i$ ساده می‌شوند:

$$\boxed{[f']_n = j\frac{2\pi n}{T}\,c_n - \frac{1}{T}\sum_i \Delta_i\,e^{-j\frac{2\pi n t_i}{T}} + \frac{1}{T}\sum_i \Delta_i\,e^{-j\frac{2\pi n t_i}{T}} = j\frac{2\pi n}{T}\,c_n}$$

پس حتی با وجود ناپیوستگی، ضرایب فوریه‌ی مشتق توزیعی $f'(t)$ دقیقاً برابر $j\dfrac{2\pi n}{T}c_n$ است. این به این دلیل است که جملات پرش (ضربه‌های $\Delta_i\delta(t-t_i)$) دقیقاً همان جملاتی را جبران می‌کنند که از انتگرال‌گیری جزء‌به‌جزء در نقاط ناپیوستگی باقی می‌مانند.


حال اگر مشتق $f'(t)$ را به صورت تابع عادی (تنها روی نقاط پیوستگی، بدون جملات ضربه) در نظر بگیریم، ضرایب فوریه‌ی آن عبارت است از:

$$\boxed{[f'_{\text{smooth}}]_n = j\frac{2\pi n}{T}\,c_n - \frac{1}{T}\sum_i \Delta_i\,e^{-j\frac{2\pi n t_i}{T}}}$$

این رابطه نشان می‌دهد که وقتی $f(t)$ پیوسته باشد (تمام $\Delta_i=0$)، رابطه‌ی $[f']_n = j\frac{2\pi n}{T}c_n$ دقیقاً برقرار است، و ناپیوستگی‌ها موجب انحراف از این رابطه می‌شوند.
